% =============================================================================
% lampiran.tex — Lampiran
% Gunakan \addlampiran{} untuk menambahkan entri ke Daftar Lampiran
% =============================================================================

% Judul Lampiran dicetak sebagai bab tanpa penomoran bab
\chapter*{LAMPIRAN}
\addcontentsline{toc}{chapter}{LAMPIRAN}
\thispagestyle{plain}

% =============================================================================
% LAMPIRAN A — Instrumen Penelitian (Kuesioner)
% =============================================================================
\section*{Lampiran A. Instrumen Penelitian}
\addlampiran{Instrumen Penelitian (Kuesioner)}
\label{lampiran:kuesioner}

\textbf{KUESIONER PENELITIAN}

\vspace{0.5cm}

\noindent
\textbf{Judul Penelitian:} \JudulTA

\vspace{0.3cm}

\noindent
Kepada Yth.\\
Bapak/Ibu Responden\\
di tempat

\vspace{0.3cm}

\noindent
Dengan hormat,

Saya yang bertanda tangan di bawah ini adalah mahasiswa Program Studi \NamaProdi,
Fakultas \NamaFakultas, \NamaUniversitas, yang sedang melakukan penelitian untuk
keperluan penyusunan tugas akhir.

Saya mohon kesediaan Bapak/Ibu untuk mengisi kuesioner ini dengan jujur sesuai
dengan keadaan yang sebenarnya. Jawaban Bapak/Ibu akan dijaga kerahasiaannya dan
hanya digunakan untuk keperluan ilmiah.

\vspace{0.5cm}

\noindent
\textbf{Petunjuk Pengisian:}
\begin{enumerate}
  \item Bacalah setiap pernyataan dengan seksama.
  \item Pilih jawaban yang paling sesuai dengan keadaan Bapak/Ibu dengan memberikan
        tanda centang ($\checkmark$) pada kolom yang tersedia.
  \item Keterangan skala: 1 = Sangat Tidak Setuju, 2 = Tidak Setuju, 3 = Netral,
        4 = Setuju, 5 = Sangat Setuju.
\end{enumerate}

\vspace{0.5cm}

\noindent\textbf{A. Identitas Responden}

\vspace{0.3cm}

\noindent
\begin{tabular}{ll}
Jenis Kelamin & : $\square$ Laki-laki \quad $\square$ Perempuan \\[0.3cm]
Usia          & : $\square$ < 25 tahun \quad $\square$ 25--35 tahun \quad $\square$ > 35 tahun \\[0.3cm]
Pendidikan    & : $\square$ SMA/SMK \quad $\square$ D3 \quad $\square$ S1 \quad $\square$ S2/S3 \\
\end{tabular}

\vspace{0.5cm}

\noindent\textbf{B. Pernyataan Penelitian}

\vspace{0.3cm}

\noindent\textit{Variabel X: [Nama Variabel]}

\begin{table}[H]
\begin{isitable}
\begin{tabular}{p{0.5cm} p{5.5cm} p{1cm} p{1cm} p{1cm} p{1cm} p{1cm}}
\toprule
\textbf{No} & \textbf{Pernyataan} & \textbf{1} & \textbf{2} & \textbf{3} & \textbf{4} & \textbf{5} \\
\midrule
1 & Pernyataan pertama variabel X & & & & & \\
2 & Pernyataan kedua variabel X & & & & & \\
3 & Pernyataan ketiga variabel X & & & & & \\
\bottomrule
\end{tabular}
\end{isitable}
\end{table}

% =============================================================================
% LAMPIRAN B — Hasil Olah Data
% =============================================================================
\clearpage
\section*{Lampiran B. Hasil Olah Data}
\addlampiran{Hasil Olah Data}
\label{lampiran:olah-data}

Lampiran ini berisi output hasil pengolahan data statistik. Tempelkan atau masukkan
output software statistik (SPSS, SmartPLS, AMOS, R, Python, dsb.) di sini.

\vspace{0.5cm}

% Contoh: output dapat disisipkan sebagai gambar
% \includegraphics[width=\textwidth]{figures/output-spss}

\noindent
[Output Uji Validitas]

\vspace{0.5cm}

\noindent
[Output Uji Reliabilitas]

\vspace{0.5cm}

\noindent
[Output Analisis Utama]

% =============================================================================
% LAMPIRAN C — Dokumen Pendukung (opsional)
% =============================================================================
\clearpage
\section*{Lampiran C. Dokumen Pendukung}
\addlampiran{Dokumen Pendukung}
\label{lampiran:dokumen}

Tambahkan dokumen pendukung penelitian di sini, misalnya:
\begin{itemize}
  \item Surat izin penelitian
  \item Panduan wawancara
  \item Transkrip wawancara
  \item Foto dokumentasi
\end{itemize}
